\section{Discussion and Conclusions}
\begin{itemize}
    \item Earthquake sequences involving partially overlapping rupture zones have been reported for several
     strike-slip fault systems, with successive large events separated by years to decades. Examples include
    the 1940 (\(M_w\) 7.1) and 1979 (\(M_w\) 6.6) earthquakes in the Imperial Valley \cite{discussion2013}, 
    the 1981 (\(M_w\) 7.1) and 1998 (\(M_w\) 6.6) earthquakes on the Gowk fault,
    Iran \cite{gowk2010}, and the 1979 (\(M_w\) 6.6), 1997 (\(M_w\) 7.2) Zirkuh earthquakes
     \cite{berberian1999} and earthquake sequences along the North Anatolian Fault \cite{barka1996}. 
     The recurrence of slip on previously ruptured sections were explained by spatial variations in co-seismic slip,
      whereas areas experiencing relatively limited slip during one
     earthquake can become sites of renewed rupture in a subsequent event. In other cases, the elapsed time between
    earthquakes was sufficient for tectonic loading to contribute substantially to the stress state
    before the subsequent rupture \cite{king1998}. 
    In contrast with the above examples, in terms of magnitude 2015 and 2025 should mark the end of a sequence of earthquake 
    cycle on the eastern segment which was previously seen in between 1974 (\(M_w\) 6.9) and 2015 (\(M_w\) 7.1) events. 
    \item How such a short gap between these two large events on the same fault segment can be explained? The main clue lies in the 
    source slip distribution of both the earthquakes. 2015 earthquake had one single confined asperity in the eastern segment of the fault
    with a maximum slip of 1.1m in the centre releasing the accumulated stress. In the other hand, 2025 earthquake had multiple small
    asperities with maximum slip of only 0.3m. From the seismic moment release point of view, to create such a large earthquake(2025)
    with such small slip amount, the rupture area must be very large which can be seen in the slip distribution (Figure~\ref{fig:2015_2025_comparison})
    \item Shear cracks with finite cohesive forces can propagate by skipping barriers, which either persist or rupture under subsequent
     dynamic stress \cite{heaton1990} depending on the barrier strength-to-tectonic stress ratio \cite{das1977}. In this context, the 2015 earthquake 
     likely left unbroken barriers and generated new ones through stress redistribution. Because the dynamic stress increase at the 
     rupture front can exceed the static stress modifications induced by the 2015 event, the 2025 rupture was able to break through 
     residual barriers while skipping newly formed ones, resulting in multiple small slip asperities.
    \item The colocation of earthquakes on the same fault system within only ten years gives clear instrumental evidence that goes
     against the idea that the ground stops shaking after a rupture. This is similar to records that show clusters of events with time
      gaps between them that are much shorter than the average cycle. The 2015 and 2025 sequence shows that fault segments that have
       already broken can break again in than a decade. Whether this happens because of a dynamic stress drop, multimodal slip or uneven 
     loading of asperities the short-interval reactivation has a direct effect on how we assess seismic hazard. A recent large
        earthquake— one that is close to the biggest expected magnitude—does not make a fault segment safe from a near‑future rupture.
     Instead new events that start in unruptured gaps or, on adjacent fault strands can spread into the area that was just ruptured.
     
\end{itemize}
